\documentclass[10pt,a4paper]{amsart}
\usepackage[margin=19mm]{geometry}
\usepackage{amsmath,amssymb,mathtools}
\usepackage[scr=rsfs,cal=euler]{mathalfa}
\usepackage{xfrac}
\usepackage[unicode,hidelinks,bookmarks=false]{hyperref}
\newcommand{\ie}{i.e.\ }
\newcommand{\cf}{cf.\ }
\newcommand{\resp}{respectively\ }
\newcommand{\Subsection}{\subsection}
\providecommand{\nobreakdash}{\nobreak\hbox{-}}
\setlength{\emergencystretch}{2em}
\multlinegap0pt
\usepackage{xcolor}
\usepackage{amssymb}
\usepackage{csquotes}
\usepackage[normalem]{ulem}
\newtheorem{proposition}{Proposition}
\title{Boundary blowup solutions for the Finsler p-Laplacian: wellposedness and asymptotic behaviour}
\author{N N Dattatreya}
\date{Source: arXiv:2605.22345v1 (CC BY 4.0)}
\begin{document}
\maketitle

\noindent\emph{Source and licence.} Boundary blowup solutions for the Finsler p-Laplacian: wellposedness and asymptotic behaviour. Version arXiv:2605.22345v1 (CC BY 4.0), distributed by arXiv under the Creative Commons Attribution 4.0 International licence, http://creativecommons.org/licenses/by/4.0/, as recorded in the arXiv licence metadata for this exact version and shown on the abstract page https://arxiv.org/abs/2605.22345v1. Full licence notice: \texttt{licenses/arxiv-2605.22345-CC-BY-4.0.txt}.\par\smallskip

\noindent\textit{Editorial scope.} Counted unit: the article's proposition proposition (source0.tex lines 561--570). The article's complete original proof of it is delivered with it (source0.tex lines 571--619, one \texttt{proof} environment of 49 lines). No mathematics is written, repaired or re-derived: the statement and the proof are the article's own lines, in the article's own notation, and no retained line is substituted or re-flowed.  Retained uncounted as the source-local context the proof needs: lines 423--428; lines 446--448; lines 467--472.  Nothing was omitted from any retained span, so the package carries no omission record.  No retained line contains a citation, so the wrapper carries no bibliography and no bibliography entry.  No retained line contains a figure environment, an image-inclusion command or drawing code, so no image is delivered and the package declares no asset dependency.  Editorial formatting only, in addition: an AMS \texttt{amsart} preamble that restates the article's own declarations and macros used by the retained lines, namely the environments \texttt{proposition} on separate counters, together with the standard packages those lines need.  The retained environments print in this document as Proposition 1 in order of appearance, whereas the article prints the same items with the numbering of its own full document.  The complete native source of the article as distributed in the arXiv e-print for this exact version is bundled unchanged as \texttt{sources/201125/source0.tex}.
        \begin{equation}\label{eqn in one dim}
        \begin{cases}
      \left(H^{p-1}(v')H'(v')\right)'=f(v)\quad \text{on } (a,b). \\
      v(x)\to \infty \quad \text{as } x\to a \text{ or }b
        \end{cases}
        \end{equation}

 \begin{equation}\label{eqn quotient}
     \left(\frac{p-1}{p}\right)^{1/p}\frac{\gamma v'(x)}{\left\{F\left(v(x)\right)-F(v_{0})\right\}^{1/p}}=1.
 \end{equation}

\begin{equation}\label{ODE2}
   \begin{cases}
     \gamma^{p}\left(|\Tilde{v}'|^{p-1}\Tilde{v}'\right)'=f(\Tilde{v})\quad \text{for all } x\leq c_{m}\\
     \Tilde{v}'(c_{m})=0 \quad \text{and} \quad \Tilde{v}(c_{m})=v_{0}.
 \end{cases} 
\end{equation}

\begin{proposition}
    Let $u$ be a solution of \eqref{eqn in one dim}. Then
    \begin{equation*}
       \lim_{x\to b}\frac{\Psi_{H,p}(u(x))}{b-x}= \frac{1}{\gamma},
    \end{equation*}
 and 
  \begin{equation*}
       \lim_{x\to a}\frac{\Psi_{H,p}(u(x))}{x-a}= \frac{1}{\gamma}.
    \end{equation*}
\end{proposition}

\begin{proof}
    Clearly 
    \begin{equation*}
        \left\{F(s)\right\}^{1/p}\geq \left\{F(s)-F(v_{0})\right\}^{1/p}.
    \end{equation*}
   On the other hand, since $F(s)\to \infty$ as $s\to\infty$, given $\epsilon>0$, there exists $s(\epsilon)>0$ such that $F(s)\geq \epsilon F(v_{0})$ for all $s\geq s(\epsilon)$;
   which implies
   \begin{equation*}
       \left\{F(s)-F(v_{0})\right\}^{1/p}\geq \left(1-\frac{1}{\epsilon}\right)^{1/p}\left\{F(s)\right\}^{1/p}
   \end{equation*}
   With these, one can write
   \begin{equation}\label{eqn double inequality}
       \left(1-\frac{1}{\epsilon}\right)^{1/p}\frac{1}{ \left\{F(s)-F(v_{0})\right\}^{1/p}}\leq \frac{1}{\left\{F(s)\right\}^{1/p}}\leq \frac{1}{ \left\{F(s)-F(v_{0})\right\}^{1/p}}.
   \end{equation}
   Since $u(t)\to\infty$ as $t\to b$, we obtain $r>0$ such that for all $t>b-r$, $u(t)\geq s(\epsilon)$. For $t>b-r$ integrating the above inequality from $u(t)$ to $\infty$, we obtain 
   \begin{equation}\label{eqn double integral ineq}
       \gamma\left(\frac{p-1}{p}\right)^{1/p}\left(1-\frac{1}{\epsilon}\right)^{1/p}\int_{u(t)}^{\infty}\frac{1}{ \left\{F(s)-F(v_{0})\right\}^{1/p}}\leq \gamma\Psi_{H,p}(u(t)) \leq \gamma\left(\frac{p-1}{p}\right)^{1/p}\int_{u(t)}^{\infty}\frac{1}{ \left\{F(s)-F(v_{0})\right\}^{1/p}}.
   \end{equation}
   Integrating \eqref{eqn quotient} from $t$ to $b$,  employing  the change of variables $s=v(x)$, and using this in the above inequality we obtain 
   \begin{equation*}
       \left(1-\frac{1}{\epsilon}\right)^{1/p}\leq \frac{\gamma\Psi_{H,p}(u(t))}{b-t}\leq 1.
   \end{equation*}
   As $t\to b$, therefore we obtain  
   \begin{equation*}
       \left(1-\frac{1}{\epsilon}\right)^{1/p}\leq \liminf_{t\to b}\frac{\gamma\Psi_{H,p}(u(t))}{b-t}\leq \limsup_{t\to b}\frac{\gamma\Psi_{H,p}(u(t))}{b-t}\leq 1.
   \end{equation*}
   and also $\epsilon\to \infty$, we obtain 
    \begin{equation*}
     \lim_{t\to b}\frac{\Psi_{H,p}(u(t))}{b-t}= \frac{1}{\gamma}.
   \end{equation*}
   %Also, note that as $\epsilon\to\infty$, $t\to b$. 
   
   Next, as $u(t)\to \infty$ as $t\to a$, there exists $r>0$ such that for all $t<a+r$ we have $u(t)>s(\epsilon)$. The counterpart of \eqref{eqn quotient} for $u=\Tilde{v}$, which solves \eqref{ODE2} is 
   \begin{equation*}
       \gamma\left(\frac{p-1}{p}\right)^{1/p}\frac{-u'(x)}{\left\{F\left(u(x)\right)-F(v_{0})\right\}^{1/p}}=1.
   \end{equation*}
   Integrating the above equation from $a$ to $t$ and using a change of variables $s=v(x)$, we obtain  
   \begin{equation*}
       t-a=\gamma\left(\frac{p-1}{p}\right)^{1/p}\int_{u(t)}^{\infty}\frac{ds}{\left\{F(s)-F(v_{0})\right\}^{1/p}}.
   \end{equation*}
   Using this equation in \eqref{eqn double integral ineq}, we have 
    \begin{equation*}
       \left(1-\frac{1}{\epsilon}\right)^{1/p}\leq \frac{\gamma\Psi_{H,p}(u(t))}{t-a}\leq 1.
   \end{equation*}
   As $\epsilon\to \infty$ we obtain 
   \begin{equation*}
       \lim_{t\to a}\frac{\Psi_{H,p}(u(t))}{t-a}=\frac{1}{\gamma}.\qedhere
   \end{equation*}
\end{proof}

\end{document}
